\subsection{Arithmetic gluing from local completions of finite form truncations}

\begin{theorem}[C34: \status{Proved under stated hypotheses}; truncated repair data suffice]
\label{sa08:C34}
Let $T$ be either of the targets $T_{\mathrm{pair}}$ or
$T_{\mathrm{eq}}$ in Statement~\ref{sa08:C32}, and let
$f:P\to T$ be its partial model from Statement~\ref{sa08:C21}.
For $d\geq1$ write
\[
 P_{<d}=P/P^{\geq d},
\]
where the ideal $P^{\geq d}$ consists of all form degrees at
least $d$. This is a quotient PD-DGA; no condition on the
outgoing differential from form degree $d-1$ survives in it.

Suppose that for each prime $p$ there are an admissible theory
$E_p$ over $k_p=k\otimes\Z_{(p)}$, a map $g_p:E_p\to T_p$,
and an integer $d_p\geq0$, with $d_p\to\infty$ as $p\to\infty$.
For each $d_p>0$, suppose there is an isomorphism of strict
simplicial weight-graded augmented PD-DGAs over $S_{k_p}$
\[
 \alpha_p:(P_p)_{<d_p}\xrightarrow{\ \cong\ }(E_p)_{<d_p}
 \quad\text{such that}\quad
 (g_p)_{<d_p}\alpha_p=(f_p)_{<d_p}.
\]
For $d_p=0$ impose no agreement condition. The isomorphisms
hold in every weight and simplex. A full map $P_p\to E_p$
is not assumed.

Then there is an admissible theory $H$ with an actual map
$H\to T$. Consequently, if these hypotheses hold for the
pair and homotopy-equalizer targets required in
Statement~\ref{sa08:C32}, the generalized category has common
sources and its localization at all arrows is terminal.
The optional condition $q\leq w$ is preserved. Over $k=\Z$,
if the targets and the $E_p$ have flat terms, $H$ may be
taken termwise flat.
\end{theorem}
\begin{proof}
Use the rings $R=\prod_p\Z_{(p)}$ and $R_\Q=R\otimes\Q$
and the pure arithmetic exact sequence proved in
Statement~\ref{sa08:C30}. Put
\[
 F_q=\{p:d_p\leq q\}\cup\{p:p\leq q+2\},
 \qquad e_{F_q}=\sum_{p\in F_q}e_p,
\]
and define
\[
 E_R^{w,q}=(1-e_{F_q})P_R^{w,q}
       \oplus\bigoplus_{p\in F_q}E_p^{w,q}.
\]
Each $F_q$ is finite. If $F\supset F_q$ is finite, every
new component has $q<d_p$, so the degree-$q$ component of
$\alpha_p$ identifies the resulting expression with the
same module. These refinements are transitive because the
fixed isomorphism $\alpha_p$ is used whenever that prime
is added.

Define a differential after refining both terms to $F_{q+1}$,
a product after refining to the set for its output form degree,
and a divided power after refining to a set containing both
its input and output degrees. On a selected prime component
use $E_p$; on the complementary idempotent use $P_R$.
Independence of refinement follows from the truncated
PD-DGA identity whenever a prime is newly added: all input
and output degrees of that operation are then strictly
below $d_p$. Every algebraic identity uses only finitely
many form degrees, so a common finite refinement proves
it componentwise. Simplicial maps, coefficients, units,
and augmentations are treated in the same way.

In particular, when $q=d_p-1$, the differential is defined
using $E_p$ after passage to $F_{q+1}$. It need not agree
with the outgoing differential of $P_p$. Likewise a product
whose output first reaches $d_p$ is the product in $E_p$.
This is precisely why a truncated isomorphism suffices and
why no unasserted map $P_R\to E_R$ is constructed.

Define $g_R:E_R\to T_R$ by $g_p$ on the selected components
and by $f_R$ on the complement. Its independence of
refinement follows from the agreement with $f_p$ in the
truncation. To check a differential or multiplicative
identity, refine to include all degrees involved and use
the corresponding identity for $g_p$ or $f_R$.

The admissibility proof of Statement~\ref{sa08:C30} now
applies without its full under-$P$ map. In each fixed row,
the complement is a flat scalar extension of the uniformly
horizontally acyclic row of $P$, and there are only finitely
many exceptional admissible rows. For vertical cohomology
through form degree $b$, refine all terms through $b+1$
to $F_{b+1}$. The complementary coefficient ring inverts
every prime at most $b+3$, so Statement~\ref{sa08:C21}
makes its $P$ complex a coefficient resolution through $b$.
The finitely many selected $E_p$ have a common eventual
weight bound in that range. Thus $E_R$ is admissible,
with bounds uniform in the simplex. Put $D=(E_R)_\Q$.

We replace the unavailable map $P_\Q\to D$ by an independent
rational elementary source. Choose a sufficiently large
power of two $L$ so that the weight-$L$, form-zero rows of
$D$, $T_\Q$, and $T_{R\Q}$ are horizontally acyclic.
Over the rational coefficient rings, the construction of
Statement~\ref{sa08:C09} defines $C=G_L$ over $k_\Q$.
Its marked lifting argument uses only form-zero horizontal
acyclicity, so it gives maps
\[
 a:C\longrightarrow T_\Q,
 \qquad r:C_{R\Q}\longrightarrow D.
\]
The latter and the rationalization of $g_R$ give two marked
maps $C_{R\Q}\rightrightarrows T_{R\Q}$, namely $a_{R\Q}$
and $g_{R\Q}r$. By the marked homotopy argument of
Statement~\ref{sa08:C16}, choose a strict simplicial PD-DGA
homotopy
\[
 K:C_{R\Q}\longrightarrow (T_{R\Q})^{\Delta[1]}
\]
from the first to the second. The proof of that marked-map
argument only uses the stated acyclicity of $T_{R\Q}^{L,0}$,
so it applies to $T_{\mathrm{eq}}$ as well.

Using $r$ after restriction of scalars, form the arithmetic
path pullback
\[
 E=C\mathop{\times}_{D,\operatorname{ev}_0}D^{\Delta[1]}
      \mathop{\times}_{\operatorname{ev}_1,D}E_R.
\]
Both endpoint theories are admissible. In a fixed eventual
range, the kernel of both endpoints of $D^{\Delta[1]}$
has only coefficient cohomology in form degree one.
The connecting map is
\[
 S_{k_\Q}^w\oplus S_{k_R}^w\longrightarrow S_{k_{R\Q}}^w,
 \qquad (u,v)\longmapsto v-u.
\]
The pure arithmetic sequence identifies its kernel with
$S_k^w$ and its cokernel with zero. The horizontal
endpoint-fibration argument and the uniform vertical
comparison in Statement~\ref{sa08:C16} therefore prove
that $E$ is admissible, exactly as in
Statement~\ref{sa08:C30}. The augmentation descends by
the same arithmetic kernel identity.

Let $I_2$ be the union of two consecutive simplicial edges,
with endpoints $0,2$, and define
\[
 J_2(T)=T_\Q\mathop{\times}_{T_{R\Q},\operatorname{ev}_0}
        (T_{R\Q})^{I_2}
        \mathop{\times}_{\operatorname{ev}_2,T_{R\Q}}T_R.
\]
For an element $(c,\gamma,e)$ of $E$, its first edge in
$J_2(T)$ is $K(c)$ and its second is $g_{R\Q}(\gamma)$.
Their common endpoint is strictly $g_{R\Q}r(c)$.
Thus these formulas give an actual map $E\to J_2(T)$;
no operation composing two simplicial homotopies into one
edge is required.

For completeness, subtracting the constant initial value
of a path gives a rowwise exact sequence
\[
 0\longrightarrow T\longrightarrow J_2(T)
   \longrightarrow P_0^{I_2}(T_{R\Q})\longrightarrow0.
\]
Surjectivity follows from the pure arithmetic sequence
applied to the final value. The based-path module is
horizontally acyclic in every row, since the pointed
two-edge interval is contractible. It is also vertically
acyclic in each eventual finite range by the
constant-cohomology relative-prism comparison of
Statement~\ref{sa08:C32}. Thus $\eta:T\to J_2(T)$ is a
horizontal weak equivalence in every row and an eventual
vertical quasi-isomorphism.

Its mapping-path factorization
\[
 \widetilde T
 =T\mathop{\times}_{J_2(T),\operatorname{ev}_0}J_2(T)^{\Delta[1]}
 \xrightarrow{\operatorname{ev}_1}J_2(T)
\]
is a trivial fibration in every simplicial module row.
The projection $\widetilde T\to T$ has an eventually
vertically acyclic based-path kernel. The two endpoint
maps on the path object induce the same vertical
cohomology map in the eventual range, so
$\widetilde T\to J_2(T)$ is also an eventual vertical
quasi-isomorphism. Its kernel is therefore horizontally
acyclic in every row and eventually vertically acyclic.
The pullback
\[
 H=E\mathop{\times}_{J_2(T)}\widetilde T
\]
is admissible and has the required actual map to $T$.
For the pair targets this gives common sources; for the
homotopy-equalizer targets it gives equality of parallel
maps after localization. The finite-zigzag argument in
Statement~\ref{sa08:C32} proves terminality.

All constructions preserve the form and weight gradings.
If $k=\Z$ and the stated terms are flat, the finite
idempotent sums are torsionfree, and the path modules
and pullbacks are subgroups of torsionfree groups.
They are therefore flat over $\Z$.
\end{proof}

\paragraph{Precise reduction achieved.}
The prescribed data at prime $p$ are now only the finite
form truncation of $P_p$ and its map to $T_p$. Higher
relations and higher values of that particular partial
model need not survive. This weakens the sufficient
hypothesis in Statement~\ref{sa08:C32}. It still asks for
a full admissible $E_p$ in every weight and simplex, and
the completion depths must still tend to infinity across
the primes. Neither existence nor necessity of these
completions has been proved. In particular this theorem
alone supplies no new coefficient case of integral C3
and does not establish that its remaining sufficient
construction is easier or equivalent to the main problem.
